% =====================================================================================================================
% Chapitre 3 : la surprise quantique (énergie de point zéro), démonstrations humaines et Lean côte à côte.
% Sources : papers/prog_zero_point_optimality, p2_certified_zero_point, p2_quantum_transition,
% prog_superfluid_structure_factor, ar_bogoliubov_sk ; Lean formal_cert/Programmes/{ZeroPoint,ZeroPointCert,FeynmanBijl}.
% Tous les nombres sont recopiés des fichiers gen_numbers.tex / gen/*.tex / main.tex de ces articles.
% =====================================================================================================================
\lstdefinestyle{chIIIlean}{language=Lean,
  literate={ℝ}{{$\mathbb{R}$}}1 {ℤ}{{$\mathbb{Z}$}}1 {ℕ}{{$\mathbb{N}$}}1 {≤}{{$\le$}}1 {≥}{{$\ge$}}1 {≠}{{$\ne$}}1
  {→}{{$\to$}}1 {↔}{{$\leftrightarrow$}}1 {∀}{{$\forall$}}1 {∃}{{$\exists$}}1 {∧}{{$\wedge$}}1 {∨}{{$\vee$}}1
  {¬}{{$\neg$}}1 {⊤}{{$\top$}}1 {∈}{{$\in$}}1 {∉}{{$\notin$}}1 {∩}{{$\cap$}}1 {⁻¹}{{$^{-1}$}}1 {√}{{$\surd$}}1
  {π}{{$\pi$}}1 {ρ}{{$\rho$}}1 {κ}{{$\kappa$}}1 {α}{{$\alpha$}}1 {δ}{{$\delta$}}1 {ε}{{$\varepsilon$}}1 {σ}{{$\sigma$}}1
  {•}{{$\bullet$}}1 {∑}{{$\sum$}}1 {²}{{$^2$}}1 {₀}{{$_0$}}1 {₁}{{$_1$}}1 {₂}{{$_2$}}1 {‖}{{$\|$}}1 {⟨}{{$\langle$}}1 {⟩}{{$\rangle$}}1
  {∞}{{$\infty$}}1 {×}{{$\times$}}1 {·}{{$\cdot$}}1 {∘}{{$\circ$}}1
  {ψ}{{$\psi$}}1 {φ}{{$\varphi$}}1 {λ}{{$\lambda$}}1 {μ}{{$\mu$}}1 {θ}{{$\theta$}}1 {τ}{{$\tau$}}1
  {—}{{---}}1 {–}{{--}}1 {é}{{\'e}}1
  {⟪}{{$\langle\!\langle$}}1 {⟫}{{$\rangle\!\rangle$}}1 {ₗ}{{$_\ell$}}1 {ℂ}{{$\mathbb{C}$}}1 {←}{{$\leftarrow$}}1}

\chapter{La surprise quantique : énergie de point zéro et réseaux}\chaptermark{La surprise quantique}\label{ch3:chap}\label{ch3:debut}

\noindent Le chapitre~1 a raconté une victoire : pour toute interaction répulsive \emph{complètement monotone}
$V(r)=g(r^2)$, le réseau triangulaire minimise l'énergie \emph{statique} $\tfrac12\sum_{R\neq0}V(|R|)$, parmi toutes les
configurations de densité un (énoncé de l'amont openai/math, \tierA{} conditionnel à des théorèmes amont produits par une
IA et non évalués par les pairs ; voir aussi le manuel version~1~\cite{ch3-txtv1}). Ce chapitre pose la question suivante, très naturelle pour un physicien : \emph{que
devient ce résultat quand on tient compte de la mécanique quantique ?} Un cristal réel n'est jamais immobile. Même au
zéro absolu, chaque atome vibre autour de son site, et cette vibration coûte une énergie, l'énergie de \emph{point zéro}.

Nous allons voir que la réponse est une surprise. Pour une large famille de gaussiennes, le réseau triangulaire n'est
pas un minimum local de l'énergie de point zéro : c'est un \emph{maximum} local (\tierX). Nous en donnerons ensuite une
version \emph{certifiée}, où la logique et l'arithmétique sont vérifiées par le noyau Lean (\tierA) et où les intégrales
sont encadrées par arithmétique de boules (\tierC). Puis nous demanderons si cela produit une vraie transition de phase
(réponse négative, dans le cadre de l'approximation harmonique auto-cohérente). Nous terminerons par un cousin de la même
mathématique : les fluides quantiques, avec la borne de Feynman--Bijl et le facteur de structure de Bogoliubov.

\begin{warnbox}[: ce que ce chapitre affirme, et ce qu'il n'affirme pas]
Tous les résultats ont été produits par des agents d'IA (Claude) sous la direction de l'auteur, et audités par un agent
d'IA séparé. \textbf{Aucun humain ne les a audités.} Rien dans ce chapitre ne dépend d'un théorème de l'amont openai/math.
Le théorème de Montgomery sur les gaussiennes, utilisé pour l'énergie statique, est cité de seconde main.
Les énoncés sur l'énergie de point zéro sont de niveau~X ou~C. Seules la logique, l'algèbre et l'arithmétique
rationnelle sont de niveau~A. L'étiquette \tierA{} s'attache toujours à l'énoncé Lean, jamais à sa lecture physique.
\end{warnbox}

%=====================================================================================================================
\section{Le cristal harmonique, pour l'élève}\label{ch3:sec:cristal}

\subsection{L'oscillateur harmonique quantique et son énergie minimale}\label{ch3:sec:oscillateur}
Une particule de masse $m$ dans le potentiel $\tfrac12m\omega^2x^2$ est décrite par une fonction d'onde
$\psi:\R\to\mathbb{C}$ de norme $\int|\psi|^2=1$. Son énergie moyenne est
\[
\langle H\rangle_\psi=\frac{\hbar^2}{2m}\int_\R|\psi'(x)|^2\,dx+\frac{m\omega^2}{2}\int_\R x^2|\psi(x)|^2\,dx .
\]
En mécanique classique, l'énergie minimale vaut $0$ (la particule au repos au fond du puits). En mécanique quantique, ce
n'est plus possible.

\begin{proposition}[Énergie de point zéro de l'oscillateur]\label{ch3:prop:osc}
Pour toute fonction $\psi$ de classe $C^1$, normée, telle que $x\psi$ et $\psi'$ soient de carré intégrable et que
$x|\psi(x)|^2\to0$ en $\pm\infty$, on a $\langle H\rangle_\psi\ge\tfrac12\hbar\omega$, avec égalité si et seulement si
$\psi(x)=c\,e^{-m\omega x^2/2\hbar}$ avec $|c|$ fixé par la normalisation. \tierC
\end{proposition}
\begin{humanproof}
Posons $\phi=m\omega\,x\psi+\hbar\psi'$. On développe $\int|\phi|^2\ge0$ :
\[
0\le\int|\phi|^2=m^2\omega^2\!\int x^2|\psi|^2+\hbar^2\!\int|\psi'|^2+2m\omega\hbar\,\mathrm{Re}\!\int x\,\overline{\psi}\,\psi' .
\]
Or $2\,\mathrm{Re}(\overline\psi\psi')=(|\psi|^2)'$, donc par intégration par parties (le terme de bord
$[x|\psi|^2]$ est nul par hypothèse)
$2\,\mathrm{Re}\int x\overline\psi\psi'=\int x(|\psi|^2)'=-\int|\psi|^2=-1$. Il reste
\[
0\le m^2\omega^2\!\int x^2|\psi|^2+\hbar^2\!\int|\psi'|^2-m\omega\hbar=2m\langle H\rangle_\psi-m\omega\hbar ,
\]
soit $\langle H\rangle_\psi\ge\hbar\omega/2$. L'égalité impose $\phi=0$, c'est-à-dire l'équation différentielle linéaire
$\psi'=-(m\omega/\hbar)x\psi$, dont les solutions sont exactement $c\,e^{-m\omega x^2/2\hbar}$.
\end{humanproof}
\noindent L'énergie $\tfrac12\hbar\omega$ est l'\emph{énergie de point zéro}. Physiquement, confiner la particule
(petit $\int x^2|\psi|^2$) oblige $\psi$ à varier vite (grand $\int|\psi'|^2$) : c'est le principe d'incertitude.

\paragraph{Plusieurs modes indépendants \tierC.} Si un système a $n$ coordonnées $q_1,\dots,q_n$ et un hamiltonien
$H=\sum_j\bigl(\tfrac{p_j^2}{2m}+\tfrac12m\omega_j^2q_j^2\bigr)$ (modes \emph{découplés}), la fonction d'onde produit
$\prod_j\psi_j(q_j)$ montre que l'énergie fondamentale est $\sum_j\tfrac12\hbar\omega_j$. Pour un cristal, la première
étape est donc de trouver les coordonnées qui découplent les vibrations : ce sont les \emph{phonons}.

\subsection{Échauffement : la chaîne linéaire}\label{ch3:sec:chaine}
Prenons $N$ masses $m$ sur un cercle, de positions d'équilibre $na$, reliées à leurs voisines par des ressorts de raideur
$K$, et notons $u_n$ le déplacement de la masse $n$ (avec $u_{n+N}=u_n$). L'énergie potentielle est
$\tfrac K2\sum_n(u_{n+1}-u_n)^2$ et l'équation de Newton s'écrit $m\ddot u_n=K(u_{n+1}+u_{n-1}-2u_n)$.

On cherche des solutions $u_n=e^{i(kna-\omega t)}$ avec $k\in\tfrac{2\pi}{Na}\Z$ (condition de périodicité). En
substituant : $-m\omega^2=K(e^{ika}+e^{-ika}-2)=-2K(1-\cos ka)=-4K\sin^2(ka/2)$, donc
\[
\omega(k)=2\sqrt{K/m}\;\bigl|\sin(ka/2)\bigr| .
\]
La transformation de Fourier discrète découple exactement les $N$ modes (elle diagonalise la matrice circulante de
l'énergie potentielle), donc l'énergie de point zéro par particule est
\[
\frac1N\sum_k\frac{\hbar\omega(k)}2\xrightarrow[N\to\infty]{}\frac{\hbar}{2}\cdot\frac1{2\pi}\int_{-\pi}^{\pi}2\sqrt{K/m}\,
|\sin(\theta/2)|\,d\theta=\frac{2}{\pi}\,\hbar\sqrt{K/m}
\]
(somme de Riemann, $\theta=ka$). Retenons trois idées. \textbf{(i)} Les fréquences sont les racines carrées des valeurs
propres d'une matrice dépendant de $k$. \textbf{(ii)} Pour $k\to0$, $\omega\approx a\sqrt{K/m}\,|k|$ : c'est le son, et
la fréquence s'annule comme $|k|$ (singularité conique). \textbf{(iii)} L'énergie de point zéro est une \emph{moyenne sur
la zone de Brillouin} $\theta\in[-\pi,\pi]$. Tout ceci est \tierC.

\subsection{Le réseau plan : matrice dynamique et énergie de point zéro}\label{ch3:sec:matdyn}
Un réseau de Bravais plan de covolume un est déterminé, à rotation près, par sa \emph{forme} $\tau=x+iy$ ($y>0$), avec la
base $a_1=(1/\sqrt y,0)$, $a_2=(x/\sqrt y,\sqrt y)$. Le carré de la longueur du vecteur $ma_1+na_2$ vaut
\[
Q_{mn}(\tau)=\frac{(m+nx)^2+n^2y^2}{y}.
\]
Le réseau triangulaire correspond à $\rho=e^{i\pi/3}=\tfrac12+i\tfrac{\sqrt3}2$, le réseau carré à $\tau=i$.

\paragraph{La hessienne d'un potentiel radial.} Pour $V(R)=g(|R|^2)$ avec $R\in\R^2$, on a
$\partial_iV=2g'(|R|^2)R_i$ puis
\begin{equation}\label{ch3:eq:hess}
\mathrm{Hess}\,V(R)=2g'(|R|^2)\,I+4g''(|R|^2)\,RR^{\mathsf T}.
\end{equation}
Pour la gaussienne $g(u)=e^{-\pi au}$ : $g'=-\pi a\,g$, $g''=\pi^2a^2g$, d'où
$\mathrm{Hess}\,V(R)=(-2\pi a\,I+4\pi^2a^2RR^{\mathsf T})e^{-\pi a|R|^2}$ \tierC.

\paragraph{Les équations du mouvement.} Plaçons une particule de masse $m=1$ en chaque point $R$ du réseau $L$, et notons
$u_R$ son petit déplacement. Au second ordre, l'énergie potentielle est
$\tfrac14\sum_{R\ne R'}(u_R-u_{R'})^{\mathsf T}\mathrm{Hess}\,V(R-R')(u_R-u_{R'})$ et Newton donne
$\ddot u_{0}=-\sum_{R\neq0}\mathrm{Hess}\,V(R)\,(u_0-u_R)$. Avec l'onde plane $u_R=e\,e^{i(\mathbf k\cdot R-\omega t)}$
($e\in\mathbb{C}^2$), on obtient $\omega^2e=\sum_{R\ne0}(1-e^{i\mathbf k\cdot R})\mathrm{Hess}\,V(R)\,e$. Les termes $R$
et $-R$ ont la même hessienne, donc on peut symétriser : $1-e^{i\mathbf k\cdot R}$ devient $1-\cos\mathbf k\cdot R$.

\begin{definition}[Matrice dynamique, stabilité, énergie de point zéro]\label{ch3:def:ezp}
Pour $\mathbf k\in\R^2$, la \emph{matrice dynamique} est la matrice symétrique réelle $2\times2$
\[
D(\mathbf k)=\sum_{R\in L\setminus\{0\}}\bigl(1-\cos\mathbf k\cdot R\bigr)\,\mathrm{Hess}\,V(R).
\]
Le réseau est \emph{dynamiquement stable} si $D(\mathbf k)$ est positive pour tout $\mathbf k$. Ses deux fréquences sont
alors $\omega_j(\mathbf k)=\sqrt{\lambda_j(D(\mathbf k))}$, et l'\emph{énergie harmonique de point zéro} par particule
(unités $\hbar=m=1$) est
\[
E_{\mathrm{zp}}(\tau)=\frac12\Bigl\langle\omega_1+\omega_2\Bigr\rangle_{\mathrm{ZB}}=\frac12\Bigl\langle
\operatorname{tr}D(\mathbf k)^{1/2}\Bigr\rangle_{\mathrm{ZB}},
\]
où $\langle\cdot\rangle_{\mathrm{ZB}}$ est la moyenne sur une maille du réseau dual (la zone de Brillouin). \tierC
\end{definition}
\noindent La formule $\operatorname{tr}D^{1/2}$ est le cœur de ce chapitre : l'énergie de point zéro est la trace de la
\emph{racine carrée} d'une somme sur le réseau. Elle n'est pas une somme de paires, contrairement à l'énergie statique.
Pour une matrice $2\times2$ positive, $\sqrt{\lambda_1}+\sqrt{\lambda_2}=\sqrt{\operatorname{tr}D+2\sqrt{\det D}}$
(élever au carré), formule utilisée par les codes numériques.

\paragraph{Une règle de somme \tierC.} La moyenne de $\cos\mathbf k\cdot R$ sur la zone est nulle pour $R\neq0$. Donc
$\langle\operatorname{tr}D\rangle_{\mathrm{ZB}}=\sum_{R\ne0}\Delta V(R)$, et par Cauchy--Schwarz puis Jensen
(exercice~\ref{ch3:ex:einstein}),
\begin{equation}\label{ch3:eq:einstein}
E_{\mathrm{zp}}\le\sqrt{\langle\operatorname{tr}D\rangle_{\mathrm{ZB}}/2}\qquad\text{(borne d'Einstein).}
\end{equation}
Pour la gaussienne, la sommation de Poisson donne
$\langle\operatorname{tr}D\rangle=4\pi a-(4\pi^2/a)\,S(1/a)$ avec $S(\alpha)=\sum_{P\in L}|P|^2e^{-\pi\alpha|P|^2}$ \tierC.
Luo et Wei démontrent que, pour $\alpha\ge1$, le minimum de $S(\alpha;\cdot)$ sur les réseaux de densité un est atteint
au réseau hexagonal~\cite{ch3-luowei} \tierL ; donc, pour $a\le1$, le triangulaire \emph{maximise}
$\langle\operatorname{tr}D\rangle$ \tierC. Une borne supérieure maximale en $\rho$ ne dit cependant rien des extrema de
$E_{\mathrm{zp}}$ lui-même.

\subsection{Le rôle du point triangulaire : symétrie et isotropie}\label{ch3:sec:isotropie}
L'énergie d'un réseau ne dépend que du réseau à rotation près ; elle est donc invariante par $\tau\mapsto\tau+1$ et
$\tau\mapsto-1/\tau$, donc par tout le groupe $SL_2(\Z)$. L'élément $\gamma(\tau)=(\tau-1)/\tau$ fixe $\rho$, car
$\rho^2=\rho-1$, et sa dérivée en $\rho$ vaut $1/\rho^2=-\tfrac12-i\tfrac{\sqrt3}2$ : c'est la rotation d'angle
$-120^\circ$. Si $F$ est de classe $C^2$ et invariante, son gradient en $\rho$ est fixé par cette rotation, donc nul, et
sa hessienne commute avec elle, donc elle est scalaire. Ainsi
\begin{equation}\label{ch3:eq:courbure}
F(\tau)-F(\rho)=c\,\frac{|\tau-\rho|^2}{y_\rho^2}+O(|\tau-\rho|^3),\qquad c=\tfrac14\Delta F(\rho),\quad
\Delta=y^2(\partial_x^2+\partial_y^2).
\end{equation}
\textbf{Le signe de $c$ décide tout :} minimum local si $c>0$, maximum local si $c<0$. Les deux faits d'algèbre
linéaire sont vérifiés par le noyau \tierA ; le passage « invariance $\Rightarrow$ gradient nul, hessienne scalaire »
(dérivation en chaîne, régularité de $E_{\mathrm{zp}}$) est \tierC.

\begin{leanbox}[: \lean{formal_cert/Programmes/ZeroPoint/Basic.lean}, lignes 91--92, 168--171 et 201]
\begin{lstlisting}[style=chIIIlean]
theorem lap_Q (m n x y : ℝ) (hy : y ≠ 0) :
    y ^ 2 * (2 * n ^ 2 / y + 2 * (m + n * x) ^ 2 / y ^ 3) = 2 * Q m n x y := by
theorem rot120_comm_symm (p q r s : ℝ) (hs : s ^ 2 = 3 / 4)
    (h11 : -1 / 2 * p + -s * q = p * (-1 / 2) + q * s)
    (h12 : -1 / 2 * q + -s * r = p * -s + q * (-1 / 2)) :
    q = 0 ∧ r = p := by
theorem rho_sq : rho ^ 2 = rho - 1 := by
\end{lstlisting}
\end{leanbox}
\noindent Lecture : \lean{lap_Q} dit que $\Delta Q=2Q$ pour la fonction longueur (on y reconnaît
$y^2(Q_{xx}+Q_{yy})$, avec $Q_{xx}=2n^2/y$ et $Q_{yy}=2(m+nx)^2/y^3$). \lean{rot120_comm_symm} dit qu'une matrice
symétrique $\bigl(\begin{smallmatrix}p&q\\q&r\end{smallmatrix}\bigr)$ commutant avec la rotation de $120^\circ$
(de coefficients $-\tfrac12$ et $\pm s$, $s^2=\tfrac34$) vérifie $q=0$ et $r=p$. Le module contient 19 déclarations
imprimées, toutes d'axiomes dans $\{$\lean{propext}, \lean{Classical.choice}, \lean{Quot.sound}$\}$, et des contrôles
piégés (un \lean{sorry} et un axiome planté) attrapés par \texttt{\#print axioms}~\cite{ch3-p2zp}. \tierA

\paragraph{L'énergie statique.} Pour la gaussienne, le théorème de Montgomery~\cite{ch3-montgomery} affirme que le réseau
triangulaire minimise l'énergie statique $E_{\mathrm{st}}=\tfrac12\sum_{R\ne0}e^{-\pi a|R|^2}$ parmi les réseaux de
covolume un, pour tout $a>0$ (\tierL, cité de seconde main à travers Deng--Luo~\cite{ch3-dengluo}). Pour un potentiel
d'amplitude $\varepsilon$, de portée $l$ et une masse $m$, l'énergie harmonique totale s'écrit
\begin{equation}\label{ch3:eq:total}
E(\tau)=E_{\mathrm{st}}(\tau)+\hbar_{\mathrm{eff}}\,E_{\mathrm{zp}}(\tau),\qquad\hbar_{\mathrm{eff}}=\frac{\hbar}{\sqrt{m\varepsilon l^2}} ,
\end{equation}
en unités de $\varepsilon$ et $l$ \tierC. Le paramètre $\hbar_{\mathrm{eff}}$ mesure « combien le système est
quantique ».

%=====================================================================================================================
\section{La conjecture QU-1 et sa réfutation numérique}\label{ch3:sec:qu1}

\subsection{L'énoncé, et la règle du jeu}
\begin{definition}[Conjecture QU-1]\label{ch3:def:qu1}
Pour toute fonction $g$ complètement monotone, le réseau triangulaire minimise $E_{\mathrm{zp}}(\tau)$ parmi les réseaux de
Bravais plans de covolume un.
\end{definition}
\noindent C'était le « compagnon quantique » naturel de l'optimalité universelle statique. Elle a été
\emph{préenregistrée} avec une probabilité a priori de $0{,}6$ de survivre à une porte de falsification numérique (G1),
une tentative de preuve pour les gaussiennes (G2) ne devant être lancée \emph{que si} G1 survivait~\cite{ch3-p2zp}.
La règle importante de la méthode Elenchus est ici : \emph{on cherche d'abord à réfuter avant de chercher à démontrer}.

\subsection{Le contre-exemple \tierX}\label{ch3:sec:contreex}
Pour la gaussienne $g(t)=e^{-\pi t}$ ($a=1$), le tableau~\ref{ch3:tab:head} compare le réseau triangulaire et le réseau
rhombique $\tau=0{,}5+0{,}966025\,i$.
\begin{table}[ht]\centering\small
\caption{Gaussienne $a=1$ : énergie de point zéro (règle de zone spectrale, erreur de quadrature
d'environ $4{,}4\times10^{-15}$), marge de stabilité $\min\lambda/|\mathbf k|^2$ sur la zone, et limite de grande
longueur d'onde exacte \tierX~\cite{ch3-p2zp}.}\label{ch3:tab:head}
\begin{tabular}{lrrrr}\toprule
réseau $\tau$ & $E_{\mathrm{zp}}$ & $\Delta E_{\mathrm{zp}}$ & $\min\lambda/|\mathbf k|^2$ & gr.\ long.\ d'onde\\\midrule
triangulaire $\rho$ & 1,5033523522 & 0 & 0,03222 & 0,23826\\
$0{,}5+0{,}966025\,i$ & 1,4902783183 & $-1{,}3074\times10^{-2}$ & 0,03493 & 0,10466\\
$i$ (carré, contrôle négatif) & -- & -- & $-0{,}1475$ (instable) & $-0{,}1475$\\
\bottomrule\end{tabular}
\end{table}

\noindent Le réseau rhombique est dynamiquement stable avec une marge nette, et son énergie de point zéro est
\emph{plus basse} que celle du réseau triangulaire. Quatre implémentations de deux formules (somme directe, et somme de
Poisson dans l'espace réciproque) concordent ; un court programme de vérification indépendant
(\texttt{human\_check.py}, dont l'article~\cite{ch3-p2zp} reproduit le cœur) redonne $1{,}490278318$ et $1{,}503352352$. Sur les 625 formes échantillonnées pour ce
potentiel, 65 sont stables et les 64 formes stables non triangulaires ont toutes une énergie plus basse : autour de
$\rho$, le triangulaire est un \emph{maximum local} dans la région stable \tierX.

\paragraph{Douze potentiels.} La conjecture est réfutée numériquement pour 8 des 12 potentiels complètement monotones
testés (tableau~\ref{ch3:tab:pot}). Le cas \texttt{riesz:4} est faible : il repose sur un seul concurrent lointain, dont
la marge de stabilité vaut 0,08\,\% de la valeur triangulaire, et le point triangulaire y est un minimum local. Les
quatre survivants ne sont \emph{pas} des preuves : un échantillon de 567 formes et deux chemins denses ne sont pas une
recherche exhaustive.
\begin{table}[ht]\centering\small
\caption{Verdicts QU-1 et nature du point triangulaire (signe de $c$, intervalle sur 12 directions) \tierX~\cite{ch3-p2zp}.
\texttt{gauss:$a$} est $e^{-\pi ar^2}$, \texttt{yukawa:$\kappa$} est $e^{-\kappa r}/r$, \texttt{riesz:$s$} est $r^{-s}$,
\texttt{dipolar:$d$} est l'exciton dipolaire.}\label{ch3:tab:pot}
\begin{tabular}{llr@{\,\ldots\,}ll}\toprule
potentiel & verdict QU-1 & \multicolumn{2}{c}{$c$ en $\rho$} & nature de $\rho$\\\midrule
\texttt{gauss:0.25} & survit & $+0{,}19$ & $+0{,}19$ & min\\
\texttt{gauss:0.5} & \textbf{réfutée} & $-0{,}47$ & $-0{,}45$ & max\\
\texttt{gauss:1} & \textbf{réfutée} & $-0{,}98$ & $-0{,}95$ & max\\
\texttt{gauss:2} & survit & $+0{,}48$ & $+0{,}51$ & min\\
\texttt{gauss:4} & survit & $+0{,}31$ & $+0{,}32$ & min\\
\texttt{yukawa:1} & \textbf{réfutée} & $-0{,}40$ & $-0{,}38$ & max\\
\texttt{yukawa:2} & \textbf{réfutée} & $-0{,}32$ & $-0{,}31$ & max\\
\texttt{dipolar:0.1} & \textbf{réfutée} & $-0{,}33$ & $-0{,}32$ & max\\
\texttt{dipolar:0.5} & \textbf{réfutée} & $-0{,}61$ & $-0{,}58$ & max\\
\texttt{riesz:3} & \textbf{réfutée} & $-0{,}31$ & $-0{,}29$ & max\\
\texttt{riesz:4} & \textbf{réfutée} (faible) & $+0{,}34$ & $+0{,}36$ & min\\
\texttt{riesz:6} & survit & $+3{,}38$ & $+3{,}45$ & min\\
\bottomrule\end{tabular}
\end{table}

\subsection{La famille gaussienne : un intervalle de maxima locaux \tierX}\label{ch3:sec:intervalle}
Pour $V=e^{-\pi ar^2}$, la courbure $c(a)=\tfrac14\Delta E_{\mathrm{zp}}(\rho)$ a été calculée pour 23 valeurs de
$a\in[0{,}1;8]$, puis ses changements de signe ont été localisés par dichotomie (tableau~\ref{ch3:tab:ca}).
\begin{table}[ht]\centering\small
\caption{Courbure $c(a)$ de $E_{\mathrm{zp}}$ au point triangulaire, gaussienne $e^{-\pi ar^2}$ (extrait) \tierX~\cite{ch3-p2zp}.}\label{ch3:tab:ca}
\setlength{\tabcolsep}{3pt}\begin{tabular}{rrrrrrrrrrr}\toprule
$a$ & 0,25 & 0,35 & 0,4 & 0,5 & 0,8 & 1 & 1,2 & 1,4 & 1,6 & 2\\\midrule
$c(a)$ & $+0{,}1882$ & $+0{,}0796$ & $-0{,}0799$ & $-0{,}4605$ & $-1{,}0707$ & $-0{,}9620$ & $-0{,}6598$ & $-0{,}3035$ & $+0{,}0285$ & $+0{,}4945$\\
\bottomrule\end{tabular}
\end{table}

\begin{proposition}[Résultat numérique, P2]\label{ch3:prop:intervalle}
Pour la gaussienne $e^{-\pi ar^2}$ avec $0{,}15\le a\le4$, la courbure $c(a)$ est négative, donc le réseau triangulaire
est un \emph{maximum local strict} de $E_{\mathrm{zp}}$, pour $a\in(0{,}37668\,;\,1{,}58148)$. \tierX
\end{proposition}
\noindent Les bornes sont des racines numériques obtenues par dichotomie ; le balayage n'exclut pas d'autres changements
de signe entre ses points, et en $c=0$ le test du second ordre ne conclut pas. La version restreinte QU-1$'$ (minimum
local strict hors de l'intervalle, maximum local strict dedans) est une conjecture soutenue au niveau~X, non démontrée.

\subsection{Le mécanisme : la concavité de la trace de la racine carrée}\label{ch3:sec:mecanisme}
Pourquoi l'analogie avec le cas statique échoue-t-elle ? L'énergie statique est \emph{linéaire} en $V$. Un potentiel
complètement monotone est un mélange positif de gaussiennes (théorème de Bernstein, chapitre~1), et une propriété vraie
pour chaque gaussienne passe au mélange. L'énergie de point zéro, elle, contient $\operatorname{tr}\sqrt{D}$ : elle n'est
pas linéaire en $V$, et l'argument de mélange ne s'applique pas \tierC.

Plus précisément, à $\mathbf k$ fixé, notons $\Phi(D)=\operatorname{tr}D^{1/2}$. Dériver deux fois
$\Phi(D(\tau))$ donne deux termes :
\[
\Delta E_{\mathrm{zp}}=\underbrace{\tfrac12\Bigl\langle y^2\textstyle\sum_i\operatorname{tr}\bigl(\Phi'(D)\,\partial_i^2D\bigr)\Bigr\rangle}_{\mathrm{Lin}\ \text{(réponse linéaire)}}
+\underbrace{\tfrac12\Bigl\langle y^2\textstyle\sum_i\Phi''(D)[\partial_iD,\partial_iD]\Bigr\rangle}_{\mathrm{Conc}\ \le\ 0} .
\]
Le second terme est négatif par \emph{concavité} de la racine carrée sur les matrices positives (l'analogue matriciel de
$\sqrt{\cdot}$ concave sur $\R_+$) \tierC. Le premier n'a pas de signe contrôlé. Pour $a=1$, la décomposition
numérique donne $\Delta E_{\mathrm{zp}}(\rho)=-3{,}848=-0{,}799+(-3{,}048)$ : la concavité domine \tierX. Pour $a=0{,}25$
et $a=2$, au contraire, un grand terme $\mathrm{Lin}>0$ l'emporte.

\begin{physbox}[: où l'énergie est-elle gagnée ?]
Pour le concurrent rhombique à $a=1$, le gain vient surtout d'un \emph{ramollissement de la branche basse près du bord
de zone}. La moyenne des carrés des fréquences baisse peu : la borne d'Einstein~\eqref{ch3:eq:einstein} n'explique que
24\,\% du gain ; les 76\,\% restants viennent de l'écart de Jensen, c'est-à-dire d'un spectre plus étalé. Comme
$\sqrt{\cdot}$ est concave, étaler le spectre à moyenne quadratique fixée fait baisser la moyenne des fréquences.
Ces nombres sont \tierX, et leur lecture est \tierC.
\end{physbox}

\begin{warnbox}[: les concurrents ne sont pas libres de contraintes]
Le réseau triangulaire a l'énergie statique la plus basse ; un concurrent rhombique n'est donc maintenu dans sa forme que
par une contrainte. QU-1 compare une fonctionnelle sur des formes imposées, et elle est réfutée comme telle. La question
physique (énergie totale, chaque réseau à sa forme d'équilibre) est celle de la section~\ref{ch3:sec:transition}.
\end{warnbox}

%=====================================================================================================================
\section{B2 : un résultat certifié}\label{ch3:sec:b2}
Un calcul en virgule flottante, même robuste, ne démontre rien. Le programme B2 a voulu un énoncé \emph{certifié} à
$a=1$~\cite{ch3-b2}. L'énoncé préenregistré demandait un encadrement bilatère certifié du laplacien $\Delta E_{\mathrm{zp}}(\rho)<0$
vérifié en Lean. \textbf{Ce n'est pas ce qui a été obtenu.} On a obtenu un énoncé plus faible sur un point (borne
unilatérale, locale, sur un seul chemin) et différent sur un autre (il n'utilise ni l'isotropie, ni la criticité).

\paragraph{Le cadre.} On part du réseau triangulaire $L_\rho$ (un vecteur sur l'axe $x$) et on le cisaille par la
matrice symétrique de déterminant un
\[
P_t=\begin{pmatrix}1&t\\t&1+t^2\end{pmatrix},\qquad L_t=P_tL_\rho,\qquad
E(t)=\frac12\int_{[-\frac12,\frac12]^2}\operatorname{tr}D_t(\kappa)^{1/2}\,d\kappa ,
\]
avec $V=e^{-\pi r^2}$ et les étiquettes du réseau et la variable $\kappa$ (coordonnée réduite) fixées quand $t$ varie.
La forme du réseau $L_t$ est $\tau(t)=\rho+\sqrt3\,t+O(t^2)$ \tierC.

\begin{theorem}[B2]\label{ch3:thm:b2}
Pour tout $t$ avec $0<|t|\le10^{-5}$, on a $E(t)<E(0)$. Par conséquent le réseau triangulaire n'est pas un minimum local de
l'énergie harmonique de point zéro pour $e^{-\pi r^2}$. \tierC{} au total ; logique et arithmétique \tierA{} ; intégrales
\tierC{} (arithmétique de boules).
\end{theorem}
\noindent La preuve tient en cinq briques. Nous les donnons dans l'ordre, chacune avec son niveau.

\subsection{Brique 1 : une inégalité de trace pour les matrices \texorpdfstring{$2\times2$}{2x2}}\label{ch3:sec:trace}
\begin{lemma}[Concavité de $\operatorname{tr}\sqrt{\cdot}$ sous forme tangente]\label{ch3:lem:trace}
Soient $X$ symétrique réelle définie positive, $S$ symétrique réelle et $P=X^{-1}$. Alors
\[
2\operatorname{tr}S\le\operatorname{tr}X+\operatorname{tr}(X^{-1}S^2),\qquad\text{avec égalité si }S=X.
\]
\tierA{} pour les matrices $2\times2$ (le seul cas utilisé) ; \tierC{} en taille quelconque.
\end{lemma}
\begin{humanproof}
\emph{Cas scalaire d'abord.} Pour $x>0$ et $s\in\R$ : $x+s^2/x-2s=(s-x)^2/x\ge0$. C'est exactement l'inégalité
arithmético-géométrique, et c'est aussi la tangente à la fonction concave $\sqrt{\cdot}$.

\emph{Cas matriciel.} Posons $M=S-X$, symétrique. En développant et en utilisant $PX=XP=I$ :
\[
\operatorname{tr}(PM^2)=\operatorname{tr}(PS^2)-\operatorname{tr}(PSX)-\operatorname{tr}(PXS)+\operatorname{tr}(PX^2)
=\operatorname{tr}(PS^2)-2\operatorname{tr}S+\operatorname{tr}X ,
\]
car $\operatorname{tr}(PSX)=\operatorname{tr}(XPS)=\operatorname{tr}S$ (invariance de la trace par permutation
circulaire), $\operatorname{tr}(PXS)=\operatorname{tr}S$ et $PX^2=X$. Il suffit donc de montrer
$\operatorname{tr}(PM^2)\ge0$. Si $m_1,m_2$ sont les colonnes de $M$, alors $M^2=MM^{\mathsf T}$ et
$\operatorname{tr}(PMM^{\mathsf T})=\operatorname{tr}(M^{\mathsf T}PM)=m_1^{\mathsf T}Pm_1+m_2^{\mathsf T}Pm_2\ge0$, puisque
$P$ est définie positive. Pour l'égalité, $S=X$ donne $M=0$.

\emph{Positivité de la forme quadratique, à la main.} Si
$P=\bigl(\begin{smallmatrix}p_{11}&p_{12}\\p_{12}&p_{22}\end{smallmatrix}\bigr)$ avec $p_{11}>0$ et $\det P\ge0$, alors
\[
p_{11}\,(p_{11}a^2+2p_{12}ab+p_{22}b^2)=(p_{11}a+p_{12}b)^2+(p_{11}p_{22}-p_{12}^2)\,b^2\ge0,
\]
et on divise par $p_{11}>0$.
\end{humanproof}

\begin{leanbox}[: \lean{formal_cert/Programmes/ZeroPointCert/Basic.lean}, lignes 49--50 et 64--78]
\begin{lstlisting}[style=chIIIlean]
theorem quad_nonneg (p11 p12 p22 a b : ℝ) (hp : 0 < p11) (hdet : 0 ≤ p11 * p22 - p12 ^ 2) :
    0 ≤ p11 * a ^ 2 + 2 * p12 * a * b + p22 * b ^ 2 := by

theorem two_trace_le_of_inverse (x11 x12 x22 p11 p12 p22 s11 s12 s22 : ℝ)
    (hp : 0 < p11) (hdet : 0 ≤ p11 * p22 - p12 ^ 2)
    (h11 : x11 * p11 + x12 * p12 = 1) (h12 : x11 * p12 + x12 * p22 = 0)
    (h21 : x12 * p11 + x22 * p12 = 0) (h22 : x12 * p12 + x22 * p22 = 1) :
    2 * (s11 + s22) ≤ (x11 + x22) +
      (p11 * (s11 ^ 2 + s12 ^ 2) + 2 * p12 * (s12 * (s11 + s22)) + p22 * (s12 ^ 2 + s22 ^ 2)) := by
  have key : (x11 + x22) +
      (p11 * (s11 ^ 2 + s12 ^ 2) + 2 * p12 * (s12 * (s11 + s22)) + p22 * (s12 ^ 2 + s22 ^ 2)) - 2 * (s11 + s22)
      = (p11 * (s11 - x11) ^ 2 + 2 * p12 * (s11 - x11) * (s12 - x12) + p22 * (s12 - x12) ^ 2)
        + (p11 * (s12 - x12) ^ 2 + 2 * p12 * (s12 - x12) * (s22 - x22) + p22 * (s22 - x22) ^ 2) := by
    linear_combination (2 * s11 - x11) * h11 + (2 * s12 - x12) * h12 + (2 * s12 - x12) * h21
      + (2 * s22 - x22) * h22
  have q1 := quad_nonneg p11 p12 p22 (s11 - x11) (s12 - x12) hp hdet
  have q2 := quad_nonneg p11 p12 p22 (s12 - x12) (s22 - x22) hp hdet
  linarith
\end{lstlisting}
\end{leanbox}
\noindent \textbf{Côte à côte.} Les quatre hypothèses \lean{h11}--\lean{h22} sont les quatre coefficients de $XP=I$. Le
membre de droite est $\operatorname{tr}X+\operatorname{tr}(PS^2)$ écrit coefficient par coefficient pour $S$ symétrique.
L'identité \lean{key} est notre « $\operatorname{tr}(PM^2)=\dots$ », écrite comme somme de deux formes quadratiques
évaluées sur les deux colonnes $(s_{11}-x_{11},s_{12}-x_{12})$ et $(s_{12}-x_{12},s_{22}-x_{22})$ de $M$ ; la tactique
\lean{linear_combination} vérifie qu'elle découle des quatre équations de $XP=I$ avec les coefficients indiqués. Les deux
appels à \lean{quad_nonneg} sont la positivité colonne par colonne. L'énoncé Lean est restreint aux matrices
$2\times2$, ce qui suffit : il y a deux branches de phonons.

\subsection{Brique 2 : le majorant de réponse linéaire \tierC}\label{ch3:sec:majorant}
En chaque $\kappa\ne0$, posons $S=D_t(\kappa)^{1/2}$ (donc $S^2=D_t$) et $X=D_0(\kappa)^{1/2}$, définie positive car le
réseau triangulaire est dynamiquement stable. Le lemme~\ref{ch3:lem:trace} donne
\[
\operatorname{tr}D_t^{1/2}\le\tfrac12\bigl(\operatorname{tr}X+\operatorname{tr}(X^{-1}D_t)\bigr).
\]
En intégrant : $E(t)\le U(t):=\tfrac14\int\bigl[\operatorname{tr}X+\operatorname{tr}(X^{-1}D_t)\bigr]d\kappa$, et
$U(0)=E(0)$ par le cas d'égalité. \textbf{Pourquoi ce majorant ?} $U$ est \emph{affine} en $D_t$ : sa dépendance en $t$
est une seule somme absolument convergente sur le réseau, contre la matrice fixe $X^{-1}$, sans racine carrée de la
matrice perturbée. Le prix est que $U$ ne borne $E$ que par au-dessus. Le terme négligé,
$\tfrac12\operatorname{tr}(X^{-1}(S-X)^2)\ge0$, ne peut qu'aider la conclusion « $E$ est plus petit ». Lean ne voit que
l'inégalité ponctuelle ; l'intégration, l'intégrabilité et l'identification $S=D_t^{1/2}$ sont de la « colle » humaine.

\subsection{Brique 3 : la borne de fenêtre}\label{ch3:sec:fenetre}
Comme $U$ est affine en $D_t$, sa différence seconde symétrique est une intégrale de la différence seconde des hessiennes
$H_n(t)=\mathrm{Hess}\,V(P_tR_n)$ :
\[
U(t)+U(-t)-2U(0)=\tfrac14\int\operatorname{tr}\Bigl(X^{-1}\textstyle\sum_n(1-\cos2\pi\kappa\cdot n)\bigl[H_n(t)+H_n(-t)-2H_n(0)\bigr]\Bigr)d\kappa .
\]
Coefficient par coefficient, la formule de Taylor--Lagrange donne $H_n(t)+H_n(-t)-2H_n(0)=t^2H_n''(\xi)$ avec
$|\xi|<t$. En séparant $H_n''(\xi)=H_n''(0)+(H_n''(\xi)-H_n''(0))$, on obtient pour $0<t\le t_0=10^{-5}$ :
\begin{equation}\label{ch3:eq:fenetre}
\frac{U(t)+U(-t)-2U(0)}{t^2}\le\underbrace{U''(0)}_{\mathrm{Upp0}}+\underbrace{\Bigl(\int\operatorname{tr}X^{-1}\Bigr)}_{\mathrm{trW}}\cdot
\underbrace{\sum_n\sum_{\text{coef.}}\sup_{|\xi|\le t_0}\bigl|H_n''(\xi)-H_n''(0)\bigr|}_{\mathrm{esum}} .
\end{equation}
La réduction est \tierC. Les dérivées le long du chemin, $q(t)=|P_tR_0|^2$, $q'$, $q''$ et la règle du produit pour
un coefficient de hessienne, sont vérifiées par le noyau (\lean{hasDerivAt_qS}, \lean{hasDerivAt_gaussS'},
\lean{hasDerivAt_entry'}) \tierA.

\paragraph{Les nombres \tierC{} (arithmétique de boules).} Les trois quantités de~\eqref{ch3:eq:fenetre} ont été encadrées
rigoureusement en arithmétique de boules Arb~\cite{ch3-arb} (python-flint 0.6.0, 96 bits, arrondi vers l'extérieur) : la
zone est découpée en quatre triangles de sommet $\Gamma$ (changement de variables de Duffy, qui supprime la singularité
conique), puis en 206\,468 cellules, sans aucun échec, en 448~s sur un cœur. Les queues de la somme sur le réseau sont
bornées explicitement (tableau~\ref{ch3:tab:boules}). Ce code n'est \emph{pas} vérifié par le noyau : une erreur dans les
bornes ne serait pas détectée par les théorèmes Lean.
\begin{table}[ht]\centering\small
\caption{Encadrements certifiés par arithmétique de boules, $a=1$, chemin de cisaillement, $t_0=10^{-5}$
\tierC~\cite{ch3-b2}.}\label{ch3:tab:boules}
\begin{tabular}{lll}\toprule
quantité & encadrement / borne & signification\\\midrule
$U''(0)=\mathrm{Upp0}$ & $[-1{,}947713\,;\,-1{,}900168]$ & dérivée seconde du majorant\\
$\int\operatorname{tr}X^{-1}\,d\kappa$ & $[1{,}825520\,;\,1{,}833534]$ & poids de la perturbation\\
$\mathrm{esum}$ & $\le0{,}0319642$ & queues comprises\\
$B_{hi}=\mathrm{Upp0}_{hi}+\mathrm{trW}_{hi}\,\mathrm{esum}_{hi}$ & $-1{,}841560$ & borne supérieure de~\eqref{ch3:eq:fenetre}\\
$\lambda_{\min}(B_0)$ & $\ge0{,}918889$ & stabilité de $L_\rho$ sur chaque cellule\\
$\|B_t-B_0\|$ & $\le0{,}007528$ & marge $0{,}911361>0$ pour $|t|\le10^{-5}$\\
\bottomrule\end{tabular}
\end{table}

\subsection{Brique 4 : le certificat rationnel}\label{ch3:sec:certificat}
Les bornes supérieures des boules sont recopiées en Lean sous forme de rationnels, et le noyau vérifie leur signe.
\begin{leanbox}[: \lean{formal_cert/Programmes/ZeroPointCert/Basic.lean}, lignes 127--144 (commentaires omis)]
\begin{lstlisting}[style=chIIIlean]
noncomputable def t0 : ℝ := 1 / 100000
noncomputable def Upp0_hi : ℝ := -475041875131 / 250000000000
noncomputable def trW_hi : ℝ := 366706720019 / 200000000000
noncomputable def esum_hi : ℝ := 31964165579 / 1000000000000
noncomputable def lmin_lo : ℝ := 183777872659 / 200000000000
noncomputable def dB_hi : ℝ := 7528036227 / 1000000000000

theorem t0_pos : 0 < t0 := by norm_num [t0]

theorem cert_B_neg : Upp0_hi + trW_hi * esum_hi < 0 := by
  norm_num [Upp0_hi, trW_hi, esum_hi]

theorem cert_stability : 0 < lmin_lo - dB_hi := by
  norm_num [lmin_lo, dB_hi]
\end{lstlisting}
\end{leanbox}
\noindent \textbf{Vérification à la main.} $\mathrm{Upp0}_{hi}=-1{,}900167500524$, $\mathrm{trW}_{hi}=1{,}833533600095$,
$\mathrm{esum}_{hi}=0{,}031964165579$ ; le produit vaut environ $0{,}0586074$ et la somme environ $-1{,}8415601<0$.
\lean{norm_num} fait le même calcul en arithmétique rationnelle exacte. L'outil LeanFlow \lean{lean.certify} a
vérifié indépendamment les deux inégalités, réécrites en entiers \tierA. Ce que le noyau \emph{ne} vérifie \emph{pas} :
que ces rationnels bornent vraiment les intégrales (c'est le contenu \tierC{} de la brique~3). Qu'ils coïncident avec le
fichier de résultats est contrôlé par un script.

\subsection{Brique 5 : la logique, puis l'assemblage}\label{ch3:sec:logique}
\begin{lemma}[Un majorant concave exclut un minimum local]\label{ch3:lem:logique}
Soient $E,U:\R\to\R$, $t_0>0$ et $B<0$. On suppose $E\le U$ sur $[-t_0,t_0]$, $U(0)=E(0)$, et
$U(t)+U(-t)-2U(0)\le t^2B$ pour $0<t\le t_0$. Alors $0$ n'est pas un minimum local de $E$. Si de plus $E$ est paire,
alors $E(t)<E(0)$ pour tout $0<t\le t_0$. \tierA
\end{lemma}
\begin{humanproof}
Pour $0<t\le t_0$ : $E(t)+E(-t)\le U(t)+U(-t)\le2U(0)+t^2B<2U(0)=2E(0)$. Donc l'un au moins des nombres $E(t)$, $E(-t)$
est strictement inférieur à $E(0)$. Si $0$ était un minimum local, il existerait $\varepsilon>0$ tel que $E\ge E(0)$ sur
$]-\varepsilon,\varepsilon[$ ; le choix $t=\min(\varepsilon/2,t_0)$ donne une contradiction. Si $E$ est paire,
$2E(t)<2E(0)$.
\end{humanproof}
\begin{leanbox}[: \lean{formal_cert/Programmes/ZeroPointCert/Basic.lean}, lignes 100--102 et 157--162]
\begin{lstlisting}[style=chIIIlean]
theorem not_isLocalMin_of_majorant {E U : ℝ → ℝ} {t0 B : ℝ} (ht0 : 0 < t0) (hB : B < 0)
    (hEU : ∀ t, |t| ≤ t0 → E t ≤ U t) (h0 : U 0 = E 0)
    (hU : ∀ t, 0 < t → t ≤ t0 → U t + U (-t) - 2 * U 0 ≤ t ^ 2 * B) : ¬ IsLocalMin E 0 := by

theorem triangular_not_localMin_of_certificate (E U : ℝ → ℝ) (Upp0 trW esum : ℝ)
    (hEU : ∀ t, |t| ≤ t0 → E t ≤ U t) (h0 : U 0 = E 0)
    (hwin : ∀ t, 0 < t → t ≤ t0 → U t + U (-t) - 2 * U 0 ≤ t ^ 2 * (Upp0 + trW * esum))
    (hU : Upp0 ≤ Upp0_hi) (hW0 : 0 ≤ trW) (hW : trW ≤ trW_hi) (he0 : 0 ≤ esum) (he : esum ≤ esum_hi)
    (hsym : ∀ t, E (-t) = E t) :
    ¬ IsLocalMin E 0 ∧ ∀ t, 0 < t → t ≤ t0 → E t < E 0 := by
\end{lstlisting}
\end{leanbox}
\noindent \textbf{Comment lire le théorème assemblé.} Chaque hypothèse est un fait que le travail analytique doit
fournir : \lean{hEU} et \lean{h0} sont la brique~2, \lean{hwin} la brique~3, \lean{hU}, \lean{hW}, \lean{he} les
encadrements de boules, \lean{hsym} la symétrie miroir $E(-t)=E(t)$ (la réflexion $\mathrm{diag}(1,-1)$ préserve $L_\rho$
et conjugue $P_t$ en $P_{-t}$, argument \tierC). Le théorème Lean est donc l'énoncé \emph{exact} de ce que les niveaux~C
doivent garantir, pas une preuve de la physique. Le titre du théorème~\ref{ch3:thm:b2} est \tierC{} au total, parce
qu'il repose sur des briques~C ; ce sont seulement sa logique et son arithmétique qui sont~A.

\paragraph{Pourquoi cela exclut un minimum local du problème des formes \tierC.} Comme $\tau(t)=\rho+\sqrt3\,t+O(t^2)$,
les petits $t\ne0$ donnent des réseaux de covolume un, de forme différente de $\rho$, arbitrairement proches de $\rho$,
et chacun a une énergie de point zéro strictement plus basse. Un minimum local en $\rho$ exigerait
$E_{\mathrm{zp}}\ge E_{\mathrm{zp}}(\rho)$ sur un voisinage, ce que contredit ce chemin. Ni la criticité ni l'isotropie
en $\rho$ ne sont utilisées.

\paragraph{Corollaire, sous hypothèses supplémentaires \tierC.} Si en outre $E$ est $C^2$, si $E_{\mathrm{zp}}$ est $C^2$
près de $\rho$, si $\rho$ est critique et la hessienne isotrope (section~\ref{ch3:sec:isotropie}), alors
$E''(0)\le-1{,}84156$ et $\Delta E_{\mathrm{zp}}(\rho)\le-0{,}92078<0$. C'est compatible avec la valeur numérique
$-3{,}848$ de P2 \tierX, et beaucoup plus faible : la méthode borne $E''$ par $U''$ d'un seul côté.

\subsection{Les contrôles}\label{ch3:sec:controlesb2}
\begin{itemize}
\item \textbf{Contrôle négatif $a=2$} (P2 y trouve un minimum local) : on obtient $U''(0)\in[+12{,}69365\,;\,+12{,}75701]$, le
certificat n'est \emph{pas} produit ; LeanFlow le rejette (\texttt{false-statement}) et le fichier \lean{RejectA2.lean}
échoue à la compilation (rc~1). Cela montre que la méthode ne dit pas toujours « pas un minimum » ; cela ne montre pas
qu'elle est complète. \tierC{} et \tierA.
\item \textbf{Certificat falsifié} (signe de $\mathrm{Upp0}_{hi}$ inversé) : \lean{RejectTampered.lean} échoue (rc~1) \tierA.
\item \textbf{\lean{sorry} et axiome plantés} : \lean{Controls.lean} compile avec rc~0, mais \texttt{\#print axioms}
affiche \lean{sorryAx} et \lean{ProgZeroPointCert.ctl_cheat}. Le code de retour n'est pas la porte ; la liste des axiomes
l'est \tierA.
\item \textbf{Sous-produit non préenregistré, $a=1/2$} : $B_{hi}=-0{,}918357<0$, même conclusion \tierC.
\end{itemize}

\begin{warnbox}[: portée exacte du niveau A ici]
Le module \lean{ZeroPointCert/Basic.lean} a 18 déclarations imprimées, toutes d'axiomes exactement
$\{$\lean{propext}, \lean{Classical.choice}, \lean{Quot.sound}$\}$ ; il n'importe jamais \lean{Mathlib} en entier.
Une recherche de \lean{skipKernelTC}, \lean{addDecl}, \lean{native_decide}, \lean{implemented_by}, \lean{sorry} et
\lean{axiom} dans les fichiers \lean{ZeroPoint/Basic.lean}, \lean{ZeroPointCert/Basic.lean} et
\lean{FeynmanBijl/Basic.lean} ne trouve rien (vérifié lors de la rédaction de ce chapitre). En revanche, aucun rejeu
\lean{leanchecker} (le troisième contrôle de la définition du niveau~A) n'est enregistré dans les dossiers de ces trois
modules (P2, B2 et P6) : il reste à faire. Leur niveau~A repose ici sur la compilation par le noyau, la liste des axiomes
et la recherche ci-dessus. Les huit drapeaux \texttt{LEDGER\_UNAUDITED\_TIER\_A} du registre restent ouverts, car aucune
personne n'a certifié que les énoncés Lean disent ce que disent leurs gloses.
\end{warnbox}

\paragraph{Travaux en cours.} Deux suites préenregistrées le 2026-10-26 sont \emph{en cours} et aucun de leurs résultats
n'est utilisé ici : N1 (un théorème de classification pour les gaussiennes : le triangulaire est-il un minimum local de
$E_{\mathrm{zp}}$ si et seulement si $a$ est hors d'un intervalle $(a_-,a_+)$ ?) et N2 (étendre le certificat de B2 aux
potentiels physiques : Yukawa, exciton dipolaire, Riesz).

%=====================================================================================================================
\section{B1 : y a-t-il une vraie transition de phase ?}\label{ch3:sec:transition}
L'énergie statique préfère le triangulaire, l'énergie de point zéro (pour $a$ dans l'intervalle) le défavorise. D'après
\eqref{ch3:eq:total}, la somme $E_{\mathrm{st}}+\hbar_{\mathrm{eff}}E_{\mathrm{zp}}$ doit donc changer de minimiseur
quand $\hbar_{\mathrm{eff}}$ grandit. Le programme B1~\cite{ch3-b1} a demandé : \emph{où, et est-ce physique ?} Sa
règle d'arrêt préenregistrée disait : « si la transition est au-delà de la fusion ou disparaît au-delà de l'harmonique,
le rapporter (résultat négatif) ». \textbf{La règle s'est déclenchée : c'est un résultat négatif.} B1 ne contient aucun
énoncé Lean.

\subsection{L'équilibre des courbures}\label{ch3:sec:equilibre}
Par~\eqref{ch3:eq:courbure} appliquée aux deux énergies, la courbure totale en $\rho$ est
$c(\hbar_{\mathrm{eff}})=c_{\mathrm{st}}+\hbar_{\mathrm{eff}}\,c_{\mathrm{zp}}$. Montgomery donne $c_{\mathrm{st}}\ge0$,
et $c_{\mathrm{zp}}<0$ sur l'intervalle de la proposition~\ref{ch3:prop:intervalle}. Le triangulaire cesse donc d'être un
minimum local (harmonique) en
\[
\hbar^*=-\frac{c_{\mathrm{st}}}{c_{\mathrm{zp}}}\qquad\text{\tierC{} (argument), \tierX{} (nombres).}
\]
À $a=1$ : $c_{\mathrm{st}}=0{,}11913$, $c_{\mathrm{zp}}=-0{,}96198$, d'où $\hbar^*=0{,}1238$. Sur $a\in[0{,}4;1{,}5]$,
$\hbar^*$ est minimal ($0{,}0877$) en $a=0{,}5$ et diverge aux deux bouts de l'intervalle, là où $c_{\mathrm{zp}}\to0$.

\paragraph{Une transition du premier ordre \tierC{} et \tierX.} Le stabilisateur de $\rho$ est le groupe diédral $D_3$.
Dans la coordonnée de disque $w=(\tau-\rho)/(\tau-\bar\rho)=re^{i\theta}$, un développement invariant contient $|w|^2$
et $\mathrm{Re}\,w^3$, mais pas $w^2$ : un \emph{terme cubique} $r^3\cos3\theta$ est permis. C'est le scénario de Landau
d'une transition du premier ordre : un état de distorsion \emph{finie} rattrape le triangulaire avant que la courbure ne
s'annule. Numériquement, à $a=1$, un réseau oblique stable $\tau\approx0{,}30+0{,}96\,i$ atteint l'énergie triangulaire dès
$\Lambda_1=0{,}0942<\hbar^*$ ; les pentes mesurées sur deux rayons miroirs sont égales et opposées (rapport $-1{,}02$),
comme le prédit le terme cubique. Les gagnants sont sur le bord de la région de stabilité harmonique.

\subsection{Le critère de Lindemann}\label{ch3:sec:lindemann}
Un cristal fond quand l'amplitude quadratique moyenne des vibrations devient une fraction notable de la distance entre
voisins : c'est la \emph{règle empirique} de Lindemann, pas un théorème. Dans la théorie harmonique à $T=0$,
$\langle|u|^2\rangle=\hbar_{\mathrm{eff}}\,U(a)$ avec $U=\tfrac12\langle\operatorname{tr}D^{-1/2}\rangle_{\mathrm{ZB}}$
(fini en dimension deux, car la singularité $1/|\mathbf k|$ est intégrable), et le rapport de Lindemann est
$\gamma=\sqrt{\langle|u|^2\rangle}/a_{nn}$ avec $a_{nn}^2=2/\sqrt3$. Comme $\gamma^2\propto\hbar_{\mathrm{eff}}$, le
couplage de fusion estimé $\Lambda_L(\gamma_c)=\gamma_c^2a_{nn}^2/U(a)$ est proportionnel à $\gamma_c^2$ : une
erreur de $\pm10\,\%$ sur $\gamma_c$ le multiplie par $0{,}81$ ou $1{,}21$ \tierC.

Les simulations Monte-Carlo quantiques de cristaux bosoniques bidimensionnels donnent, mot à mot,
« $\gamma=0.230(6)$ » pour des dipôles, « $0.235(15)$ for Yukawa bosons », « $0.24(1)$ for Coulomb bosons » et
« a hard-disks $\gamma=0.279$ »~\cite{ch3-astra} (\tierC{}, citation vérifiée mot à mot ; la porte de citation
n'admet au niveau~L que des théorèmes). D'où la bande $\gamma_c\in[0{,}23\,;\,0{,}28]$.

\paragraph{Lecture \tierX.} Au point spinodal $\hbar^*$, $\gamma\ge0{,}2798$ pour tous les $a$ testés ($0{,}313$ à
$a=1$) : au sommet de la bande ou au-delà. Au point du premier ordre $\Lambda_1$, $\gamma$ vaut $0{,}245$ à $0{,}273$ pour
$a=0{,}5$ à $1$ : \emph{dans} la bande. Le test de Lindemann seul est donc \textbf{non concluant}.

\subsection{L'approximation harmonique auto-cohérente (SCHA)}\label{ch3:sec:scha}
On remplace la matrice dynamique nue par une matrice effective $D_{\mathrm{eff}}$, déterminée variationnellement, de sorte
que le potentiel soit moyenné sur les fluctuations quantiques elles-mêmes. Le déplacement relatif $w_R=u_R-u_0$ est
alors gaussien, de covariance $C_R=\hbar_{\mathrm{eff}}\langle(1-\cos\mathbf k\cdot R)D_{\mathrm{eff}}^{-1/2}\rangle$.
Pour un potentiel gaussien, tout est explicite : avec $b=\pi a$, en complétant le carré dans l'intégrale gaussienne,
\begin{equation}\label{ch3:eq:lissage}
\bigl\langle e^{-b|R+w|^2}\bigr\rangle_{w\sim\mathcal N(0,C)}=\det(I+2bC)^{-1/2}\exp\!\bigl(-b\,R^{\mathsf T}(I+2bC)^{-1}R\bigr)
\qquad\tierC
\end{equation}
(exercice~\ref{ch3:ex:gauss}). Les fluctuations \emph{lissent} le potentiel : c'est l'effet anharmonique principal.

\paragraph{Résultat \tierX.} Dans la SCHA, à $a=1$ et $\hbar_{\mathrm{eff}}=0{,}15$, le laplacien $\Delta E(\rho)$ vaut
$+0{,}117$ (minimum local), alors que la théorie harmonique donne $\Delta E(\rho)=4(c_{\mathrm{st}}+\hbar_{\mathrm{eff}}c_{\mathrm{zp}})=-0{,}101$
(plus un minimum). Le triangulaire reste un minimum local et le
plus bas des réseaux stables échantillonnés jusqu'à $\hbar_{\mathrm{eff}}\approx0{,}27$--$0{,}33$ ; là, le rapport de
Lindemann SCHA vaut $0{,}60$ à $0{,}87$, bien au-delà de toute estimation de fusion. Deux implémentations (Python et le
module Rust \texttt{scha.rs} de rusty-SUNDIALS, commit \texttt{9e85e15}) concordent à $8{,}9\times10^{-16}$ ; un auditeur
d'IA séparé a reproduit le verdict avec son propre code et une autre quadrature.

\begin{warnbox}[: ce que dit exactement le résultat négatif]
La transition \emph{disparaît dans la SCHA}. Ce n'est pas une preuve d'absence physique : la SCHA est une théorie
variationnelle gaussienne qui omet le terme cubique du second ordre, et le durcissement qu'elle trouve est lui-même de cet
ordre. Le Monte-Carlo par intégrales de chemin (PIMC) n'a \textbf{pas} été lancé : \textbf{BLOQUÉ} (il faudrait des
mouvements de forme de boîte à aire fixe, et la taille finie masque le signal jusqu'à $M\gtrsim12$). Seuls des réseaux de
Bravais ont été comparés. L'interprétation « la fusion quantique ordinaire passe avant » est une inférence à partir de
nombres~X, pas un énoncé de la littérature.
\end{warnbox}

\paragraph{Les atomes de Rydberg habillés ne sont pas gaussiens \tierC{} et \tierX.} L'interaction d'atomes habillés par
un état de Rydberg a la forme d'un « épaulement mou » $U_0/(1+(r/R_c)^6)$~\cite{ch3-zeiher,ch3-henkel}. Sa transformée de
Fourier plane a des valeurs négatives (minimum relatif $-0{,}0448$ en $kR_c=4{,}82$, calcul \tierX), alors que celle de la
gaussienne $e^{-\pi ar^2}$ est $a^{-1}e^{-\pi k^2/a}>0$. Henkel, Nath et Pohl écrivent qu'à cause du plateau intérieur
($r<R_c$) le potentiel dans l'espace des moments a des contributions négatives autour de $k_{\mathrm{rot}}\sim2\pi/R_c$, et
ils trouvent pour ce potentiel un ramollissement de type roton et la formation d'un supersolide~\cite{ch3-henkel}.
Prestipino et al.\ écrivent, en résumant la littérature classique sur les amas, que la répulsion gaussienne est « a
marginal case »~\cite{ch3-prestipino}. Ces deux citations ont été vérifiées mot à mot. L'idée, plus générale, que les
cristaux en amas exigent une composante de Fourier négative n'est rapportée que par une lecture secondaire non revérifiée
et ne porte aucune affirmation ici. L'habillage de Rydberg ne réalise donc pas le modèle gaussien de ce chapitre.

%=====================================================================================================================
\section{Fluides quantiques : hélium superfluide et condensats}\label{ch3:sec:fluides}
Un superfluide n'est pas un cristal, mais la même mathématique y apparaît : une énergie d'excitation, une racine carrée,
et une inégalité entre une quantité mesurable (le facteur de structure statique $S(k)$, obtenu par diffraction) et le
spectre.

\subsection{Le facteur de structure de Bogoliubov et son enveloppe}\label{ch3:sec:bog}
Pour un gaz de Bose dilué de masse $m$ et de vitesse du son $c$, posons $h_c=\hbar c$, $k_*=2mc/\hbar$ et
$\varepsilon_0(k)=\hbar^2k^2/2m=h_ck^2/k_*$. La relation de dispersion de Bogoliubov et le facteur de structure
correspondant sont
\[
\varepsilon_{\mathrm B}(k)=\sqrt{(h_ck)^2+(h_ck^2/k_*)^2},\qquad S_{\mathrm B}(k)=\frac{\varepsilon_0(k)}{\varepsilon_{\mathrm B}(k)} .
\]
En Lean, $S_{\mathrm B}$ est une \emph{définition} ; la relation $\varepsilon_0/S_{\mathrm B}=\varepsilon_{\mathrm B}$ est
donc une identité algébrique, pas une dérivation à $N$ corps indépendante.

\begin{proposition}[Enveloppe de Bogoliubov]\label{ch3:prop:enveloppe}
Pour $h_c,k,k_*>0$ : $S_{\mathrm B}(k)^2=k^2/(k_*^2+k^2)$ et $S_{\mathrm B}(k)\le\sqrt{k/2k_*}$, avec égalité si et
seulement si $k=k_*$. \tierA
\end{proposition}
\begin{humanproof}
$S_{\mathrm B}^2=\varepsilon_0^2/\varepsilon_{\mathrm B}^2=\dfrac{h_c^2k^4/k_*^2}{h_c^2k^2+h_c^2k^4/k_*^2}=\dfrac{k^2}{k_*^2+k^2}$
($h_c$ disparaît). Les deux membres étant positifs,
$S_{\mathrm B}^2\le k/2k_*\iff2k_*k^2\le k(k_*^2+k^2)\iff0\le k(k-k_*)^2$, ce qui est vrai, avec égalité si et seulement
si $k=k_*$. On conclut en prenant les racines carrées, croissantes sur $[0,+\infty[$.
\end{humanproof}
\begin{leanbox}[: \lean{formal_cert/Programmes/FeynmanBijl/Basic.lean}, lignes 262--268 (commentaires omis) et 314--319]
\begin{lstlisting}[style=chIIIlean]
noncomputable def epsFree (hc ks k : ℝ) : ℝ := hc * k ^ 2 / ks
noncomputable def epsBog (hc ks k : ℝ) : ℝ := Real.sqrt ((hc * k) ^ 2 + (hc * k ^ 2 / ks) ^ 2)
noncomputable def sBog (hc ks k : ℝ) : ℝ := epsFree hc ks k / epsBog hc ks k

theorem sBog_le_sqrt {hc ks k : ℝ} (hhc : 0 < hc) (hk : 0 < k) (hks : 0 < ks) :
    sBog hc ks k ≤ Real.sqrt (k / (2 * ks)) := by
  rw [← Real.sqrt_sq (sBog_pos hhc hk hks).le]
  apply Real.sqrt_le_sqrt
  rw [sBog_sq hhc hk hks, div_le_div_iff₀ (by positivity) (by positivity)]
  nlinarith [mul_nonneg hk.le (sq_nonneg (k - ks))]
\end{lstlisting}
\end{leanbox}
\noindent \textbf{Côte à côte.} La première ligne de preuve écrit $S_{\mathrm B}=\sqrt{S_{\mathrm B}^2}$, la deuxième
réduit à comparer les carrés, la troisième remplace $S_{\mathrm B}^2$ par sa forme close et chasse les dénominateurs, et
la dernière est notre « $0\le k(k-k_*)^2$ », fournie comme indice à \lean{nlinarith}. Ce résultat a été démontré d'abord
sur une autre pile Lean~\cite{ch3-bsk}, puis redémontré sur la pile figée de ce livre~\cite{ch3-fb}.

\subsection{La borne de Feynman--Bijl, en algèbre linéaire}\label{ch3:sec:fb}
L'argument de Feynman et Bijl~\cite{ch3-feynman} (ces articles originaux n'ont pas été lus ; nous suivons la version
formalisée de~\cite{ch3-fb}) borne l'énergie de la plus basse excitation de densité par $\hbar^2k^2/(2mS(k))$.
\paragraph{Le cadre fini.} Soit $E=\mathbb{C}^n$ muni du produit hermitien $\langle x,y\rangle$ (linéaire en $y$), $H$
une matrice hermitienne, $\psi$ un vecteur propre, $H\psi=E_0\psi$ avec $E_0$ réel, et $\rho$ une matrice quelconque,
d'adjointe $\rho^\dagger$. On note $[A,B]=AB-BA$ et on définit l'\emph{excès}
\[
\mathcal A(X)=\langle X\psi,HX\psi\rangle-E_0\langle X\psi,X\psi\rangle=\langle X\psi,(H-E_0)X\psi\rangle .
\]
En physique, $\rho=\hat\rho_{\mathbf k}=\sum_je^{-i\mathbf k\cdot\mathbf r_j}$, $\psi$ est l'état fondamental et
$NS(k)=\|\rho\psi\|^2$.

\begin{lemma}[Double commutateur]\label{ch3:lem:dc}
Pour toute matrice $\rho$ : $\langle\psi,[\rho^\dagger,[H,\rho]]\psi\rangle=\mathcal A(\rho)+\mathcal A(\rho^\dagger)$. \tierA
\end{lemma}
\begin{humanproof}
On développe $[\rho^\dagger,[H,\rho]]=\rho^\dagger H\rho-\rho^\dagger\rho H-H\rho\rho^\dagger+\rho H\rho^\dagger$ et on
prend chaque terme dans $\psi$.
(a) $\langle\psi,\rho^\dagger H\rho\psi\rangle=\langle\rho\psi,H\rho\psi\rangle$ (définition de l'adjoint).
(b) $\langle\psi,\rho^\dagger\rho H\psi\rangle=E_0\|\rho\psi\|^2$ car $H\psi=E_0\psi$.
(c) $\langle\psi,H\rho\rho^\dagger\psi\rangle=\langle H\psi,\rho\rho^\dagger\psi\rangle=E_0\|\rho^\dagger\psi\|^2$
($H$ hermitienne, $E_0$ réel).
(d) $\langle\psi,\rho H\rho^\dagger\psi\rangle=\langle\rho^\dagger\psi,H\rho^\dagger\psi\rangle$.
En regroupant (a)$-$(b) et (d)$-$(c), on trouve $\mathcal A(\rho)+\mathcal A(\rho^\dagger)$.
\end{humanproof}

\begin{lemma}[Quotient de Rayleigh dans un secteur]\label{ch3:lem:rayleigh}
Soit $(e_i)$ une base orthonormée de vecteurs propres de $H$, $He_i=\lambda_ie_i$. Pour tout $x\ne0$, il existe $i$ avec
$\langle e_i,x\rangle\neq0$ et $\lambda_i\le\langle x,Hx\rangle/\|x\|^2$. \tierA
\end{lemma}
\begin{humanproof}
Avec $c_i=\langle e_i,x\rangle$ : $\langle x,Hx\rangle=\sum_i|c_i|^2\lambda_i$ et $\|x\|^2=\sum_i|c_i|^2>0$ (Parseval). Le
quotient est donc une moyenne des $\lambda_i$ pondérée par $w_i=|c_i|^2\ge0$. Si tous les $\lambda_i$ de poids $w_i>0$
étaient strictement au-dessus de la moyenne $\bar\lambda$, on aurait $\sum w_i\lambda_i>\bar\lambda\sum w_i=\sum
w_i\lambda_i$, contradiction.
\end{humanproof}

\begin{theorem}[Borne de Feynman--Bijl, forme abstraite]\label{ch3:thm:fb}
Supposons $N,m>0$, $S=\|\rho\psi\|^2/N>0$, la règle de somme f
$\langle\psi,[\rho^\dagger,[H,\rho]]\psi\rangle=N\hbar^2k^2/m$ et la symétrie $\mathcal A(\rho^\dagger)=\mathcal
A(\rho)$. Alors il existe un vecteur propre $e_i$ avec $\langle e_i,\rho\psi\rangle\ne0$ et
\[
\lambda_i-E_0\le\frac{\hbar^2k^2}{2mS}. \tierA
\]
\end{theorem}
\begin{humanproof}
Par le lemme~\ref{ch3:lem:dc} et la symétrie, $2\mathcal A(\rho)=N\hbar^2k^2/m$. L'état d'essai
$\phi=\rho\psi/\|\rho\psi\|$ a donc l'énergie
$\langle\phi,H\phi\rangle=E_0+\mathcal A(\rho)/\|\rho\psi\|^2=E_0+\hbar^2k^2/(2mS)$. Le lemme~\ref{ch3:lem:rayleigh}
appliqué à $x=\rho\psi$ fournit $e_i$.
\end{humanproof}
\begin{leanbox}[: \lean{formal_cert/Programmes/FeynmanBijl/Basic.lean}, lignes 45--46, 57--58 et 189--194]
\begin{lstlisting}[style=chIIIlean]
noncomputable def excess (H : E →ₗ[ℂ] E) (E0 : ℝ) (ψ : E) (X : E →ₗ[ℂ] E) : ℂ :=
  ⟪X ψ, H (X ψ)⟫_ℂ - (E0 : ℂ) * ⟪X ψ, X ψ⟫_ℂ

theorem double_comm_expect (hH : H.IsSymmetric) (hψ : H ψ = (E0 : ℂ) • ψ) (ρ : E →ₗ[ℂ] E) :
    ⟪ψ, comm (adjoint ρ) (comm H ρ) ψ⟫_ℂ = excess H E0 ψ ρ + excess H E0 ψ (adjoint ρ) := by

theorem feynman_bijl_bound (hH : H.IsSymmetric) (hψ : H ψ = (E0 : ℂ) • ψ) (ρ : E →ₗ[ℂ] E)
    {N m hbar k S : ℝ} (hN : 0 < N) (hm : 0 < m) (hS : 0 < S) (hSdef : S = ‖ρ ψ‖ ^ 2 / N)
    (hf : ⟪ψ, comm (adjoint ρ) (comm H ρ) ψ⟫_ℂ = ((N * hbar ^ 2 * k ^ 2 / m : ℝ) : ℂ))
    (hsym : excess H E0 ψ (adjoint ρ) = excess H E0 ψ ρ) :
    ∃ i, ⟪hH.eigenvectorBasis rfl i, ρ ψ⟫_ℂ ≠ 0 ∧
      hH.eigenvalues rfl i - E0 ≤ hbar ^ 2 * k ^ 2 / (2 * m * S) := by
\end{lstlisting}
\end{leanbox}
\noindent Ici \texttt{E} $\to_\ell[\mathbb{C}]$ \texttt{E} désigne les applications $\mathbb{C}$-linéaires de $E$ (espace de dimension finie) et
$\langle\!\langle\cdot,\cdot\rangle\!\rangle_{\mathbb{C}}$ le produit hermitien. Les quatre lignes \lean{h1}--\lean{h4}
de la preuve Lean de \lean{double_comm_expect} sont exactement (a)--(d). Le module de Feynman--Bijl compte 21
déclarations dont les axiomes sont dans la liste autorisée ; il a été recompilé indépendamment par un auditeur d'IA séparé,
qui a retrouvé les mêmes listes d'axiomes et attrapé les contrôles plantés ainsi que deux mutations de son cru~\cite{ch3-fb}.

\paragraph{D'où vient la règle de somme f ? \tierC.} Pour $N$ bosons de hamiltonien
$H=\sum_jp_j^2/2m+V(\mathbf r_1,\dots,\mathbf r_N)$, seul le terme cinétique ne commute pas avec $\hat\rho_{\mathbf k}$, et
le calcul avec $[x,p]=i\hbar$ donne l'identité d'opérateurs $[\hat\rho_{\mathbf k}^\dagger,[H,\hat\rho_{\mathbf k}]]=N\hbar^2k^2/m$.
Or la relation $[x,p]=i\hbar$ \emph{n'a aucune réalisation en dimension finie} : la trace d'un commutateur est nulle,
alors que celle de $i\hbar I$ ne l'est pas (exercice~\ref{ch3:ex:trace}). C'est pourquoi le théorème Lean, posé en
dimension finie, prend la règle de somme f en \emph{hypothèse}. Il en va de même de la symétrie $\mathbf k\leftrightarrow-\mathbf k$.
Un contre-exemple vérifié par le noyau montre que la symétrie ne peut pas être retirée de l'étape de division par deux
\tierA, et un témoin vérifié par le noyau montre que toutes les hypothèses sont satisfaisables ensemble et que la borne est
atteinte \tierA.

\paragraph{Lien avec Bogoliubov.} Pour $S>0$ : $\varepsilon_0/S\le\varepsilon_{\mathrm B}\iff S_{\mathrm B}\le S$, et
$\varepsilon_0/S_{\mathrm B}=\varepsilon_{\mathrm B}$ (\lean{feynman_le_epsBog_iff}, \lean{feynman_sBog_eq}) \tierA.
D'où le \emph{test de diffraction unilatéral} \tierC : si un $S(k)$ mesuré dépasse $\sqrt{k/2k_*}$, il n'est pas le
$S_{\mathrm B}$ d'un gaz de Bogoliubov homogène de même vitesse du son. Un point sous l'enveloppe ne confirme rien.

\subsection{Le test sur données tabulées : la région des rotons}\label{ch3:sec:roton}
\paragraph{Données \tierL.} Le tableau~14.1 de la compilation de Donnelly et Barenghi~\cite{ch3-db} (édition web,
compilation \emph{secondaire}, mesures entre 1,4 et 4~K, sans barres d'erreur imprimées) donne 99 points $(Q,S(Q))$ après
retrait de $Q=0$ (49 rayons~X, 50 neutrons). Avec $c=238{,}3$~m/s~\cite{ch3-godfrin}, on trouve
$k_*=2m_4c/\hbar=3{,}00379$~\AA$^{-1}$, donc $2k_*=6{,}008$~\AA$^{-1}$.

\paragraph{Arithmétique \tierX.} Avec un seuil $\Delta_{ds}=0{,}082$ fixé \emph{a posteriori} (l'écart entre les maxima des
deux jeux de données ; il n'est pas dans le préenregistrement), 63 points sont réfutés, 29 compatibles et 7 indécis. Dans
la région des rotons $1{,}7\le Q\le2{,}4$~\AA$^{-1}$, \textbf{les 21 points sont réfutés}, avec $S$ entre $1{,}62$ et
$2{,}49$ fois la borne. Ce résultat ne dépend pas du seuil : la plus petite marge y vaut $0{,}328$, environ $4\Delta_{ds}$.

\paragraph{Vérification à la main.} Au pic, $Q=2{,}058$~\AA$^{-1}$ et $S=1{,}458$, alors que
$\sqrt{Q/2k_*}=\sqrt{2{,}058/6{,}008}\approx0{,}585$. La marge vaut environ $0{,}87$, bien au-delà de l'écart $0{,}082$
entre les jeux de données.

\paragraph{Lecture \tierC.} L'hélium~4 liquide n'est pas un gaz de Bogoliubov avec sa propre vitesse du son. C'est un fait
de manuel (l'hélium n'est pas dilué) et \emph{pas} une découverte. La lecture est conditionnelle aux données secondaires,
à des températures de 1,4 à 4~K prises pour l'état fondamental, et au seuil a posteriori. Pour les gaz de Bose piégés,
seules des valeurs lues sur des figures existent (\tierX) : les trois points absolus de Stamper-Kurn et al.~\cite{ch3-sk}
sont sous une borne moyennée en densité locale, avec une faible puissance de test, et les points de Steinhauer et
al.~\cite{ch3-st} ont été multipliés par 2,3 par les auteurs. Aucun test décisif n'est possible.

%=====================================================================================================================
\section{Bilan du chapitre}\label{ch3:sec:bilan}
\begin{table}[ht]\centering\small
\caption{Les affirmations du chapitre et leur niveau.}\label{ch3:tab:bilan}
\begin{tabular}{p{8.6cm}p{5.6cm}}\toprule
affirmation & niveau\\\midrule
$E_{\mathrm{zp}}\ge\tfrac12\hbar\omega$ pour l'oscillateur ; phonons de la chaîne & \tierC\\
identités de $Q$, isotropie (algèbre $2\times2$), $\rho^2=\rho-1$ & \tierA{} (colle : \tierC)\\
QU-1 réfutée pour 8 potentiels sur 12 ; maximum local pour $a\in(0{,}37668\,;1{,}58148)$ & \tierX\\
mécanisme : concavité de $\operatorname{tr}\sqrt{\cdot}$ & \tierC{} (nombres \tierX)\\
inégalité $2\operatorname{tr}S\le\operatorname{tr}X+\operatorname{tr}(X^{-1}S^2)$, logique, certificat rationnel & \tierA\\
encadrements des intégrales ($U''(0)\in[-1{,}948\,;-1{,}900]$) & \tierC{} (arithmétique de boules)\\
B2 : $\rho$ n'est pas un minimum local à $a=1$ & \tierC{} au total\\
B1 : pas de transition dans la SCHA ; PIMC bloqué & \tierX{} (négatif)\\
$S_{\mathrm B}\le\sqrt{k/2k_*}$ ; borne de Feynman--Bijl abstraite & \tierA{} (règle f : hypothèse)\\
$^4$He : 21/21 points de la région des rotons au-dessus de l'enveloppe & données \tierL, calcul \tierX, lecture \tierC\\
\bottomrule\end{tabular}
\end{table}
\noindent La leçon méthodologique est double. D'abord, \emph{tester avant de démontrer} : la conjecture QU-1, plausible,
aurait pu absorber des mois de tentatives de preuve. Ensuite, \emph{séparer les étages} : un énoncé comme le
théorème~\ref{ch3:thm:b2} mélange un noyau vérifié (logique, arithmétique), du calcul rigoureux mais non vérifié par le
noyau (boules) et des arguments humains ; le théorème Lean assemblé rend cette séparation explicite, hypothèse par
hypothèse.

%=====================================================================================================================
\section{Exercices}\label{ch3:sec:exercices}

\begin{exercise}[Oscillateur]\label{ch3:ex:osc}
(a) Vérifier que $\psi_0(x)=e^{-m\omega x^2/2\hbar}$ satisfait $-\tfrac{\hbar^2}{2m}\psi_0''+\tfrac12m\omega^2x^2\psi_0=\tfrac12\hbar\omega\,\psi_0$.
(b) Dans la démonstration de la proposition~\ref{ch3:prop:osc}, justifier l'intégration par parties sur $[-A,A]$ puis le
passage à la limite. (c) Pour l'état d'essai $\psi_\alpha(x)\propto e^{-\alpha x^2}$, calculer $\langle H\rangle$ en
fonction de $\alpha$ et retrouver le minimum.
\end{exercise}

\begin{exercise}[Chaîne linéaire]\label{ch3:ex:chaine}
(a) Montrer que la matrice $N\times N$ de l'énergie potentielle de la chaîne est circulante, et que les vecteurs
$(e^{ikna})_n$, $k\in\tfrac{2\pi}{Na}\Z$, la diagonalisent. (b) Établir
$\frac1N\sum_k\tfrac12\hbar\omega(k)\to\tfrac2\pi\hbar\sqrt{K/m}$. (c) Calculer la vitesse du son $\lim_{k\to0}\omega/|k|$.
\end{exercise}

\begin{exercise}[Hessienne radiale]\label{ch3:ex:hess}
(a) Démontrer la formule~\eqref{ch3:eq:hess}. (b) En déduire $\Delta V(R)=4(ug')'(u)$ en $u=|R|^2$, puis, pour
$g(u)=e^{-\pi au}$, $\Delta V=(-4\pi a+4\pi^2a^2u)e^{-\pi au}$ (c'est la fonction \lean{m2f} de
\lean{ZeroPoint/Basic.lean}). (c) Pourquoi $D(\mathbf k)$ est-elle symétrique réelle ?
\end{exercise}

\begin{exercise}[Borne d'Einstein]\label{ch3:ex:einstein}
(a) Montrer que $\tfrac12(\omega_1+\omega_2)\le\sqrt{(\omega_1^2+\omega_2^2)/2}$. (b) En utilisant l'inégalité de
Cauchy--Schwarz pour l'intégrale sur la zone, en déduire~\eqref{ch3:eq:einstein}. (c) Cas d'égalité ? Expliquer pourquoi
une borne supérieure maximale en $\rho$ ne localise pas pour autant les extrema de $E_{\mathrm{zp}}$.
\end{exercise}

\begin{exercise}[Rotation de $120^\circ$]\label{ch3:ex:rot}
Soit $R=\bigl(\begin{smallmatrix}-1/2&-s\\s&-1/2\end{smallmatrix}\bigr)$ avec $s^2=3/4$. (a) Montrer qu'un vecteur
$g$ tel que $Rg=g$ est nul. (b) Montrer qu'une matrice symétrique $H$ telle que $RH=HR$ est scalaire (comparer avec
l'énoncé \lean{rot120_comm_symm}). (c) Montrer que $\rho^2=\rho-1$ et que $(\tau-1)/\tau$ fixe $\rho$.
\end{exercise}

\begin{exercise}[L'inégalité de trace]\label{ch3:ex:trace2}
(a) Démontrer le lemme~\ref{ch3:lem:trace} quand $X$ et $S$ sont diagonales. (b) Donner un exemple de matrices $X$, $S$
symétriques, $X$ définie positive, pour lequel l'inégalité est stricte, et vérifier numériquement. (c) Montrer que, pour
$X$ fixée, $S\mapsto\operatorname{tr}X+\operatorname{tr}(X^{-1}S^2)-2\operatorname{tr}S$ atteint son minimum $0$ en
$S=X$. Interpréter comme « la tangente est au-dessus d'une fonction concave ».
\end{exercise}

\begin{exercise}[La logique du majorant]\label{ch3:ex:logique}
(a) Rédiger complètement la preuve du lemme~\ref{ch3:lem:logique}. (b) Construire $E$ et $U$ avec $E\le U$, $U(0)=E(0)$
et $U''(0)>0$, mais où $0$ \emph{n'est pas} un minimum local de $E$. En déduire que la méthode de B2 ne peut jamais
certifier un minimum local. (c) Pourquoi l'hypothèse $E\le U$ ne suffit-elle pas si l'on remplace $U(0)=E(0)$ par
$U(0)\ge E(0)$ ?
\end{exercise}

\begin{exercise}[Le certificat rationnel]\label{ch3:ex:cert}
(a) Avec les rationnels de la section~\ref{ch3:sec:certificat}, calculer $\mathrm{Upp0}_{hi}+\mathrm{trW}_{hi}\,\mathrm{esum}_{hi}$
à $10^{-6}$ près. (b) Jusqu'à quelle valeur de $\mathrm{esum}_{hi}$ le certificat resterait-il valable ? (Réponse :
environ $1{,}0363$.) (c) Démontrer l'énoncé de monotonie \lean{cert_window} : si $\mathrm{Upp0}\le\mathrm{Upp0}_{hi}$,
$0\le\mathrm{trW}\le\mathrm{trW}_{hi}$ et $0\le\mathrm{esum}\le\mathrm{esum}_{hi}$, alors
$\mathrm{Upp0}+\mathrm{trW}\,\mathrm{esum}\le\mathrm{Upp0}_{hi}+\mathrm{trW}_{hi}\,\mathrm{esum}_{hi}$. Où sert
l'hypothèse $0\le\mathrm{esum}$ ?
\end{exercise}

\begin{exercise}[Intégrales gaussiennes et SCHA]\label{ch3:ex:gauss}
(a) En dimension un, pour $w\sim\mathcal N(0,\sigma^2)$, montrer que
$\langle e^{-b(R+w)^2}\rangle=(1+2b\sigma^2)^{-1/2}\exp\bigl(-bR^2/(1+2b\sigma^2)\bigr)$ en complétant le carré.
(b) En déduire~\eqref{ch3:eq:lissage} pour une covariance $C=\sigma^2I_2$, puis pour $C$ symétrique définie positive
quelconque en diagonalisant $C$. (c) Montrer que la transformée de Fourier plane de $e^{-\pi ar^2}$ est
$a^{-1}e^{-\pi k^2/a}$ (avec la convention $\hat f(k)=\int f(x)e^{-2i\pi k\cdot x}dx$), donc strictement positive.
\end{exercise}

\begin{exercise}[Équilibre des courbures et Lindemann]\label{ch3:ex:lind}
(a) Avec $c_{\mathrm{st}}=0{,}11913$ et $c_{\mathrm{zp}}=-0{,}96198$ ($a=1$), retrouver $\hbar^*\approx0{,}1238$.
(b) Montrer que $\Lambda_L(\gamma_c)\propto\gamma_c^2$ et que $\Lambda_1/\Lambda_L=(\gamma(\Lambda_1)/\gamma_c)^2$.
(c) Avec $\gamma(\Lambda_1)=0{,}273$ à $a=1$, dire pour quelles valeurs de $\gamma_c$ dans $[0{,}23\,;\,0{,}28]$ le point du
premier ordre précède la fusion estimée. Conclure sur le caractère non concluant du test.
\end{exercise}

\begin{exercise}[Enveloppe de Bogoliubov]\label{ch3:ex:bog}
(a) Démontrer la proposition~\ref{ch3:prop:enveloppe}. (b) Montrer le régime des phonons :
$\frac{k}{k_*+k}<S_{\mathrm B}(k)<\frac{k}{k_*}$. (c) Montrer que $S_{\mathrm B}(k)\to1$ quand $k\to\infty$ et que la valeur
libre $S\equiv1$ viole l'enveloppe pour tout $0<k<2k_*$ : c'est le contrôle positif du test. (d) Refaire la vérification
du pic de l'hélium~4.
\end{exercise}

\begin{exercise}[Pourquoi la règle f est une hypothèse]\label{ch3:ex:trace}
(a) Montrer que pour deux matrices carrées $A,B$ de taille $n$, $\operatorname{tr}(AB-BA)=0$. En déduire qu'il n'existe
pas de matrices $X,P$ avec $XP-PX=i\hbar I_n$. (b) Dans $\mathbb{C}^2$ de base orthonormée $(e_0,e_1)$, soit
$H=|e_1\rangle\langle e_1|$, $\psi=e_0$ et $\rho=|e_1\rangle\langle e_0|$. Calculer $\mathcal A(\rho)$,
$\mathcal A(\rho^\dagger)$ et le double commutateur, et vérifier que $\mathcal A(\rho)\ne f/2$ : la symétrie est
nécessaire. (c) Avec $\rho=|e_1\rangle\langle e_0|+|e_0\rangle\langle e_1|$, $N=\hbar=k=1$ et $m=1/2$, vérifier que la
borne du théorème~\ref{ch3:thm:fb} est atteinte.
\end{exercise}

\label{ch3:fin}
